\section*{الفصل \ref{ch:g12:work-and-energy} --- الشغل والطاقة الميكانيكية}

\begin{solution}{exo:g12:work-and-energy:1}
$E_p = \tfrac12 \times 200 \times 0.050^2 = \qty{0.25}{J}$. ومضاعفة
الاستطالة: $\times 4$، أي\ \qty{1.0}{J}. أما \qty{5.0}{cm} \omterm{def:g10:orders-of-magnitude:unit}{الثانية}
وحدها فتكلّف $1.0 - 0.25 = \qty{0.75}{J}$ --- إذ تكون \omterm{def:g10:inertia:force}{القوة} هناك أكبر.
\end{solution}

\begin{solution}{exo:g12:work-and-energy:2}
$W = mg(z_A - z_B) = 0.90 \times 9.81 \times (12 - 30) \approx
\qty{-1.6e2}{J}$، مهما يكن الطريق --- وهو نفسه على الصعود
المستقيم. وفي رحلة الذهاب والإياب: \qty{0}{J} (\omterm{def:g12:work-and-energy:conservative}{فالقوة محافظة}).
\end{solution}

\begin{solution}{exo:g12:work-and-energy:3}
(أ) و(ب) محافظتان: \omterm{def:g11:work-of-force:work}{فالشغل} تثبّته الطرفان (هبوط الارتفاع؛
والاستطالة). و(ج) و(ه) غير محافظتين: إذ يفرض $-fL$ وتفرض \omterm{def:g11:circuits-and-power:resistance}{المقاومة} رسمًا على كل \omterm{def:g10:orders-of-magnitude:unit}{متر}
مقطوع. و(د) عديمة \omterm{def:g11:work-of-force:work}{الشغل}: فهي عمودية دائمًا على الحركة.
\end{solution}

\begin{solution}{exo:g12:work-and-energy:4}
$h = L(1 - \cos\qty{60}{\degree}) = \qty{1.0}{m}$؛
$v = \sqrt{2 \times 9.81 \times 1.0} \approx \qty{4.4}{m/s}$. أما
\omterm{def:g10:inertia:inventory}{الشدّ} فعمودي على الحركة في كل لحظة: فلا \omterm{def:g11:work-of-force:power}{استطاعة} له،
ولا شغل.
\end{solution}

\begin{solution}{exo:g12:work-and-energy:5}
$F = P/v = 250/9.0 \approx \qty{28}{N}$. وعند \qty{5.0}{m/s}:
$F = 250/5.0 = \qty{50}{N}$ (وهي في معظمها مركّبة الثقل على المنحدر).
\end{solution}

\begin{solution}{exo:g12:work-and-energy:6}
$E_p = \tfrac12 \times 450 \times 0.080^2 = \qty{1.44}{J}$؛
$v = \sqrt{2 \times 1.44/0.020} = \qty{12}{m/s}$؛
$h = E_p/mg = 1.44/(0.020 \times 9.81) \approx \qty{7.3}{m}$.
\end{solution}

\begin{solution}{exo:g12:work-and-energy:7}
أشباه المنحرفات (\qty{2.0}{cm} خطوات): $(0.03 + 0.10 + 0.19 + 0.30) =
\qty{0.62}{J}$. والنموذج المستقيم: $k = 18.0/0.080 =
\qty{225}{N/m}$، $\tfrac12 kx^2 = \qty{0.72}{J}$. فالمنحني المقيس
يترهّل تحت المستقيم عند $x$ الصغيرة، ومن ثم تكون المساحة \omterm{def:g11:lenses-and-eye:images}{الحقيقية}
أصغر.
\end{solution}

\begin{solution}{exo:g12:work-and-energy:8}
$\Delta E_m = \tfrac12 \times 55 \times 9.0^2 - 55 \times 9.81 \times
8.0 = 2228 - 4316 \approx \qty{-2.1e3}{J}$؛
$f = 2089/25 \approx \qty{84}{N}$. \omterm{def:g10:universal-gravitation:weight}{فالوزن} محافظ، \omterm{def:g12:newtons-laws:friction}{والقوة الناظمية}
عديمة \omterm{def:g11:work-of-force:work}{الشغل}: ولا يدخل إلا الهبوط و\emph{طول} \omterm{def:g10:relative-motion:trajectory}{المسار}، لا
شكله.
\end{solution}

\begin{solution}{exo:g12:work-and-energy:9}
$H_{\min} = \tfrac52 R = \qty{0.50}{m}$. ومن \qty{60}{cm}:
$v_{\text{القمة}} = \sqrt{2 \times 9.81 \times (0.60 - 0.40)} \approx
\qty{2.0}{m/s}$، و$v_{\text{القمة}}^2/(gR) = 2.0$ --- أي ضعف
الحدّ الأدنى. \omterm{def:g10:inertia:inventory}{والاحتكاك} يقشط \omterm{def:g10:energy-conservation:energy}{الطاقة} على طول الطريق، ومن ثم تحتاج المضامير \omterm{def:g11:lenses-and-eye:images}{الحقيقية} ارتفاعًا
إضافيًّا.
\end{solution}

\begin{solution}{exo:g12:work-and-energy:10}
نقطتا الارتداد: $8.0\,x^2 = 2.0$، $x = \pm\qty{0.50}{m}$. \omterm{def:g12:kinematics-2d:velocity}{والسرعة}
القصوى عند $x = 0$: $v = \sqrt{2 \times 2.0/0.100} \approx
\qty{6.3}{m/s}$. وهي تذبذب ذهابًا وإيابًا في وادٍ مكافئ ---
و$E_p$ الخاصة \omterm{prop:g12:oscillators-and-time:spring-period}{بمذبذب الزنبرك والكتلة} من هذه الصورة بالضبط.
\end{solution}

\begin{solution}{exo:g12:work-and-energy:11}
$\tfrac12 kx^2 = \tfrac12 mv^2$:
$x = v\sqrt{m/k} = 2.0\sqrt{1200/\num{6.0e5}} \approx \qty{8.9}{cm}$.
وعند \qty{4.0}{m/s}: \qty{18}{cm} --- أي الضعف، لا أربعة أضعاف: \omterm{def:g10:energy-conservation:energy}{فالطاقة}
تتربّع، لكن $E_p$ ينمو مثل $x^2$، ومن ثم لا يتضاعف $x$ إلا مرة واحدة.
\end{solution}

\begin{solution}{exo:g12:work-and-energy:12}
(أ) $E_m = \qty{5.0}{J} < \qty{6.0}{J}$: فلا تقدر أن تعبر؛ بل تقف على
الجانب حيث $E_p = \qty{5.0}{J}$ وتتدحرج عائدة. (ب) وعند قمة التلّ:
$E_k = \qty{1.0}{J}$، $v = \sqrt{2 \times 1.0/0.50} = \qty{2.0}{m/s}$؛
وفي الوادي: $E_k = \qty{5.0}{J}$، $v \approx \qty{4.5}{m/s}$. (ج) و$x = 0$:
\omterm{prop:g12:work-and-energy:equilibria}{غير مستقر} (عظمى)؛ و$x = \qty{1.5}{m}$: \omterm{prop:g12:work-and-energy:equilibria}{مستقر} (صغرى).
\end{solution}

\begin{solution}{exo:g12:work-and-energy:13}
$\tfrac12 mv^2 = mgR_E$:
$v = \sqrt{2 \times 9.81 \times \num{6.37e6}} \approx \qty{1.12e4}{m/s}
= \qty{11.2}{km/s}$. وكل حدّ يتناسب مع $m$: فالمسبار
والحصاة سواء. ودون ذلك بقليل، ما يزال الإطلاق يصعد بعيدًا جدًّا،
ويلاقي \omterm{def:g12:work-and-energy:diagram}{نقطة ارتداد}، ويسقط عائدًا: أي مقيَّدًا.
\end{solution}

\begin{solution}{exo:g12:work-and-energy:14}
$E_m = \tfrac12 kA^2 = \tfrac12 \times 40 \times 0.060^2 =
\qty{7.2e-2}{J}$؛ $v_{\max} = \sqrt{2E_m/m} =
\sqrt{2 \times 0.072/0.250} \approx \qty{0.76}{m/s}$؛
$E_k = E_p$ حيث $\tfrac12 kx^2 = E_m/2$: $x = \pm A/\sqrt 2 \approx
\pm\qty{4.2}{cm}$.
\end{solution}

\begin{solution}{exo:g12:work-and-energy:15}
$h = v^2/2g = 9.5^2/19.62 \approx \qty{4.6}{m}$؛ والعارضة $\approx
4.6 + 1.0 = \qty{5.6}{m}$. والرقم القياسي أعلى: إذ يضيف القافزون شغلًا
بالذراعين والساقين على الزانة المنحنية. لكن تقدير \omterm{def:g10:energy-conservation:energy}{الطاقة} يضبط
الرتبة --- فلن تبلغ قفزة عدّاء \qty{10}{m} أبدًا.
\end{solution}

\begin{solution}{pb:g12:work-and-energy:1}
\textbf{1.} $v = \sqrt{2gH} = \sqrt{2 \times 9.81 \times 45} \approx
\qty{29.7}{m/s}$: \omterm{def:g10:universal-gravitation:weight}{فالوزن} محافظ، \omterm{def:g12:newtons-laws:friction}{والقوة الناظمية} عديمة \omterm{def:g11:work-of-force:work}{الشغل} --- ويسقط
المنحني من الحساب.

\textbf{2.} البوابة: $E_m = mgH = 70 \times 9.81 \times 45 \approx
\qty{3.09e4}{J}$. والانطلاق: $E_m = \tfrac12 \times 70 \times 26^2
\approx \qty{2.37e4}{J}$.

\textbf{3.} $\Delta E_m \approx \qty{-7.2e3}{J}$؛ و
مبرهنة \omterm{def:g11:mechanical-energy:em}{الطاقة الميكانيكية} تحمّلها للقوى \omterm{def:g12:work-and-energy:conservative}{غير المحافظة}:
أي \omterm{def:g10:inertia:inventory}{احتكاك} الثلج \omterm{def:g11:circuits-and-power:resistance}{ومقاومة} الهواء.

\textbf{4.} $f = \num{7.2e3}/90 \approx \qty{80}{N}$.

\textbf{5.} $\num{7.2e3}/\num{3.09e4} \approx 23\%$. والتشميع يخفض $f$،
والانحناء يخفض \omterm{def:g11:circuits-and-power:resistance}{المقاومة}: فكلاهما يقلّص $fL$، وهو الحدّ الوحيد القابل للتفاوض.

\textbf{6.} \omterm{def:g12:projectile-motion:freefall}{سقوط حرّ} \omterm{def:g12:kinematics-2d:velocity}{بسرعة} ابتدائية أفقية --- أي حركة
\omterm{def:g12:projectile-motion:projectile}{مقذوف}. $v_x = \qty{26}{m/s}$؛
$v_y = \sqrt{2 \times 9.81 \times 40} \approx \qty{28.0}{m/s}$.

\textbf{7.} $v = \sqrt{26^2 + 28.0^2} = \sqrt{1461} \approx
\qty{38.2}{m/s}$؛ وبالطاقة: $v = \sqrt{v_0^2 + 2gh} = \sqrt{676 + 785}$
--- أي متطابقان. ونحو \qty{138}{km/h}.

\textbf{8.} $t = v_y/g = 28.0/9.81 \approx \qty{2.9}{s}$؛
$d = v_0 t \approx \qty{74}{m}$.

\textbf{9.} لا. فالرفع عمودي على \omterm{def:g12:kinematics-2d:velocity}{السرعة} (عديم \omterm{def:g11:work-of-force:work}{الشغل})
\omterm{def:g11:circuits-and-power:resistance}{والمقاومة} تعاكسها (شغل سالب): $\Delta E_m \leq 0$، ومن ثم لا يمكن أن تكون \omterm{def:g12:kinematics-2d:velocity}{سرعة}
الهبوط \omterm{def:g11:lenses-and-eye:images}{الحقيقية} إلا \emph{أصغر} --- فالهواء يمطّ
الطيران، ولا يغذّيه.

\textbf{10.} $\alpha = \arctan(28.0/26) \approx \qty{47}{\degree}$
تحت الأفقي.

\textbf{11.} الزاوية مع المنحدر: $47 - 35 = \qty{12}{\degree}$.
$v_\parallel = 38.2\cos\qty{12}{\degree} \approx \qty{37.4}{m/s}$؛
$v_\perp = 38.2\sin\qty{12}{\degree} \approx \qty{7.9}{m/s}$.

\textbf{12.} الممتصّ: $\tfrac12 \times 70 \times 7.9^2 \approx
\qty{2.2e3}{J}$، في مقابل $\tfrac12 \times 70 \times 38.2^2 \approx
\qty{5.1e4}{J}$ إجمالًا: أي نحو $4\%$.

\textbf{13.} $h_{\text{المكافئ}} = 7.9^2/19.62 \approx \qty{3.2}{m}$ ---
أي قفزة جريئة من جدار. وعلى الأرض المستوية: $38.2^2/19.62 \approx \qty{74}{m}$ ---
وهي غير قابلة للنجاة.

\textbf{14.} المنحدر الموازي \omterm{def:g10:relative-motion:trajectory}{لمسار} الطيران يبقي $v_\perp$
صغيرة، ومن ثم تمتصّ الساقان أمتارًا، لا جبل \omterm{def:g11:mechanical-energy:kinetic}{الطاقة الحركية}
كله.

\textbf{15.} $H'_0 = v_0'^2/2g = 576/19.62 \approx \qty{29.4}{m}$.

\textbf{16.} $mgH' - f\,\dfrac{H'}{\sin\qty{35}{\degree}} =
\tfrac12 m v_0'^2$.

\textbf{17.} $H' = \dfrac{\tfrac12 \times 70 \times 576}
{686.7 - 80/0.574} = \dfrac{\num{20160}}{547} \approx \qty{37}{m}$.

\textbf{18.} $36.8/29.4 \approx 1.25$: أي زيادة قدرها $25\%$.

\textbf{19.} من أجل \qty{60}{kg}: $H' = \num{17280}/(588.6 - 139.5)
\approx \qty{38.5}{m}$. فالقافزة \emph{الأخف} تحتاج البرج
الأعلى: إذ تكون فاتورة \omterm{def:g10:inertia:inventory}{الاحتكاك} $fL$ هي نفسها، لكن ميزانيتها الطاقية
$mgH'$ أصغر.

\textbf{20.} أعلِن برجًا قدره \qty{37}{m} من أجل \qty{24}{m/s} عند
الانطلاق --- واعترف بأن ربعه، أي نحو \qty{7}{m} من
الخرسانة، هو زيادة \omterm{def:g10:inertia:inventory}{الاحتكاك} على الحلم الخالي من \omterm{def:g10:inertia:inventory}{الاحتكاك}.
\end{solution}
